\section{基于 Scaling Laws 的工程决策流程}

\begin{figure}[H]
\centering
\includegraphics[width=0.95\textwidth]{slides-images/slide-030.jpg}
\caption{Scaling Law 驱动的工程决策流程\protect\footnotemark}
\end{figure}
\footnotetext{Slides 第31页。}

综合本讲所有内容，基于 Scaling Laws 的模型开发流程可以总结为：

\begin{importantbox}{Scaling Law 驱动的设计流程}
\begin{enumerate}
    \item \textbf{训练小模型}：跨越几个数量级的计算量，训练一组小模型
    \item \textbf{建立 Scaling Law}：验证 log-log 线性关系确实成立
    \item \textbf{基于预测设置超参数}：利用 Scaling Law 的外推，确定大模型的最优配置
\end{enumerate}
\end{importantbox}

在许多情况下，Scaling 曲线对不同超参数的斜率相同，仅有常数偏移。这意味着你甚至不需要拟合完整的 Scaling Law——只需在\textbf{任意一个}小尺度上选出最优超参数，结论就能直接迁移。

但有一个重要例外：\textbf{学习率}需要特殊处理（通过 MuP 等方法或通过专门的学习率 Scaling Law）。

\subsection{本章小结}

Scaling Law 驱动的工程流程将大模型开发从``试错''转变为``预测+验证''。虽然并非所有决策都能完美迁移（特别是学习率和下游任务性能），但这套方法论大幅降低了探索成本。

%% ========== 总结与延伸 ==========

